\section{Introduction}
\label{sec:intro}

Mechanistic differential-equation models of signalling are useful in proportion
to what their parameters mean. In a model of the colonic crypt, a rate constant
for APC-dependent degradation of \bcat{} is not a fitting coefficient: it is the
quantity that separates a healthy crypt from an adenoma, and estimating it from
data is the point of the
exercise~\citep{fearon1990genetic,kinzler1996hereditary,clevers2012wnt}.
Two machine-learning ideas have changed how such estimates are attempted.
Physics-informed neural networks (PINNs) recover parameters by training a
network to fit the data while satisfying the differential
equations~\citep{raissi2019pinn,karniadakis2021piml}, and have been applied to
systems-biology inverse problems from sparse, noisy
measurements~\citep{yazdani2020systems,lagergren2020binns,daneker2023identifiability,ahmadi2024aiaristotle}.
Universal differential equations and neural--mechanistic hybrids keep the
trusted parts of a model and learn the untrusted part, typically one regulatory
law, with a small embedded
network~\citep{rackauckas2020ude,philipps2025ude,giampiccolo2024hnode,derooij2025ude}.

Both literatures report parameter recovery the same way: a table of estimates,
a count of how many fall within some tolerance of the truth on synthetic data,
and, for the ones that do not, an attribution. The attribution is usually one of
two kinds. Either the parameter is ``non-identifiable''~\citep{raue2009identifiability,wieland2021identifiability},
or the estimator struggled: PINN training is known to be stiff, to balance its
loss terms badly, and to stall~\citep{wang2021gradient,wang2022ntk,krishnapriyan2021failure},
and in hybrid models the over-parameterised network component is reported to
have ``consistently hindered parameter identifiability''~\citep{philipps2025ude}. The two
attributions call for opposite remedies. An estimator failure is fixed by a
better estimator; an information failure is fixed only by a better experiment,
and no amount of estimator engineering will touch it. Yet the two are routinely
diagnosed by inspection, and an improvement in the count from a new architecture
is indistinguishable, in such a table, from an improvement that a new
experimental condition would have given for free.

This paper argues that the question \emph{estimator or experiment?} has a
quantitative answer that costs one sensitivity computation, and that asking it systematically
changes the conclusions one draws. The answer is the Cram\'er--Rao lower bound
(CRLB)~\citep{rao1945information,cramer1946mathematical} of the experiment
actually performed: the Fisher information of the design at the true parameters,
inverted. It says, before any estimator is trained, how precisely every
parameter \emph{can} be known from these data. It is local and asymptotic, so we
verify it empirically; but where it holds it splits every recovery failure into
two parts that we call the \emph{estimator gap}, the distance between what an
estimator achieves and what the bound permits, and the \emph{information gap},
the parameters the bound itself declares unknowable. The first is closed by
numerics, the second only by design~\citep{kreutz2009design,bandara2009oed,apgar2010sloppy}.

We carry this accounting through a single model at increasing levels of
difficulty, from the forward problem to the inverse problem to the hybrid one.
The model is a reduced seven-state description of the WNT, retinoic acid (RA)
and HOX antagonism in the colonic crypt with 36 kinetic parameters, studied in
four regimes from a healthy crypt to severe APC loss and observed under a
ten-condition design of retinoid and WNT/MYC perturbations (Fig~\ref{fig:design}).
Four findings organise the paper.

\begin{enumerate}[leftmargin=1.4em,itemsep=0.3ex,topsep=0.4ex]
\item \emph{The PINN's recovery shortfall is an estimator gap.} The best of
three matched PINN variants recovers 21, 19, 10 and 10 of the 34 regularly
behaved parameters in the four regimes, a count that reads like a hard limit of
the data. The bound says otherwise: an efficient estimator would recover 132.5
of the 136 regime--parameter pairs outside one structural degeneracy, a
converged maximum-likelihood fit attains the bound's standard deviations on
independent datasets where the data are informative and never does worse, and
on the benchmark data, given a few starts, it recovers every one of them.
\item \emph{The estimator gap of a PINN can be closed without leaving the PINN
framework.} Freezing the network's hidden features and optimising only its
output layer together with the parameters turns the inverse problem into a
separable least-squares problem whose Gauss--Newton matrix is block-arrowhead.
A Schur-complement Gauss--Newton PINN (SGN-PINN) built on this matches the
classical estimator parameter by parameter wherever its continuation converges,
which on two datasets includes both low-WNT regimes and, on the held-out one,
Advanced Adenoma, and it reports an estimate of its own Cram\'er--Rao precision
as a by-product; where the continuation stalls, its physics residual says so.
\item \emph{The parameter that defines disease stage is the one the design does
not inform.} APC functionality $\thP$ enters only through a product with a
degradation constant; with that structural degeneracy removed, its bound still
widens from 0.04\% to 57\% relative standard deviation (that of the identifiable
product it enters, thirty-fold) as WNT drive saturates \bcat{}. No estimator can recover what the bound withholds; one added WNT
knockdown restores it, and two added conditions that perturb without leaving
the saturated plateau do not.
\item \emph{A learned regulatory term creates an information gap of its own.}
When a flexible network replaces a production term $a+f(u)$, the constraint
$f(0)=0$ that is meant to separate the basal rate $a$ from the learned function
carries information only if the experiment drives the regulator $u$ near zero,
which no condition of a conventional design does. The basal rate is then lost
under every network parameterisation we tried. A Fisher-information analysis of
the hybrid equation, consistent with fits run before it, shows that two added
conditions that do not visit the anchor leave the loss unchanged while a
depletion knockdown that visits it removes most of the loss; and a
training-free design table built on the idea predicted the right knockdown
prospectively, including one informative failure that the Fisher analysis
explains.
\end{enumerate}

The thread through these findings is that a recovery count, by itself, does not
say whether to change the estimator or the experiment, and the bound does. The
same yardstick that exonerates the data in the first two findings
indicts them in the last two, and in both cases it names the remedy: a solver
where the gap is numerical, a perturbation that visits the region in which the
parameter acts where the gap is informational. We frame the latter as a
partial-persistence-of-excitation
condition~\citep{wang2006learning,sastry1989adaptive} applied to a model
constraint rather than to an estimate. All data are synthetic, so that truth
is known and every claim can be checked; the limitations this imposes are
discussed in Section~\ref{sec:discussion}.
